% sigma(z), the smooth stand-in for step(z) that the logistic classifier uses.
%
% Used on lecture 4 badges 13 and 28, which share one raster. Deliberately the
% same drawing as step.tex: same axes, same range, same ticks, same line weight,
% same label placement. The two plots sit side by side on badge 13, so any
% difference in styling reads as a difference in meaning.
%
% Three things differ from step.tex, all because the curve differs:
%   - the guide at 1 runs the full width, since sigma approaches 1 from both
%     sides rather than reaching it at a threshold
%   - there is a 0.5 tick, the value the classifier compares against
%   - the curve is sampled, not three straight segments
%
% Build:  pdflatex sigmoid.tex  &&  ../pdf2png.sh sigmoid

\documentclass[tikz]{standalone}
\usepackage{pgfplots}
\usepackage{amsmath}
\pgfplotsset{compat=1.18}

\begin{document}
\begin{tikzpicture}
  \begin{axis}[
    axis lines=middle,
    x axis line style={->,>=stealth, shorten >=-3pt},
    y axis line style={->,>=stealth, shorten >=-3pt},
    xmin=-4.4, xmax=4.4,
    % ymin at 0 so the y axis has no tail below the origin, as in step.tex
    ymin=0, ymax=1.42,
    xtick={-4,-2,2,4},
    ytick={0.5,1},
    yticklabel style={anchor=east},
    extra x ticks={0}, extra x tick style={xticklabel style={anchor=north east}},
    tick style={black, thin},
    axis line style={black},
    width=10cm, height=4.75cm,
    scale only axis,
    clip=false,
  ]
    % the asymptote, drawn before the curve so the curve sits on top
    \addplot[gray, dashed, thin, domain=-4:4, samples=2] {1};

    % sigma(z) = 1 / (1 + e^{-z})
    \addplot[black, line width=2pt, domain=-4:4, samples=200] {1/(1+exp(-x))};

    % Axis names as nodes rather than xlabel/ylabel: both sit beside their
    % arrow tip, which pgfplots' label anchoring will not do on a middle axis.
    \node[anchor=south west] at (axis cs:0.12,1.20) {\(\sigma(z)\)};
    \node[anchor=south east] at (axis cs:4.38,0.05) {\(z\)};
  \end{axis}
  \pgfresetboundingbox
  \useasboundingbox ([xshift=-9mm, yshift=-7mm]current axis.south west)
    rectangle ([xshift=6mm, yshift=6mm]current axis.north east);
\end{tikzpicture}
\end{document}
